% Nouveau EGA V (cote 25): archive cote 25, site page 181 (PDF page 182).
% English translation of the checked French annotated master.
% AI-translated by GPT-6 Astra (OpenAI, ChatGPT Pro mode).
\scanpage{25/181}
\mspage{\hand{8}}

of a direct sum of forms $Q_1 \tikz[baseline=(sum.base)]{\node[draw,circle,handwritten,text=black,inner sep=0pt] (sum) {$\unsure{+}$};} Q_2$ , $Q_1$ of rank 1 \unsure{and} $Q_2$ of even rank.
By 16.8. we see that $Q_1,Q_2$ have corrected discriminant
$\hand{\neq} 0$; by the preceding discussion it follows that $Q_2$ is isomorphic to the
standard form, and on the other hand $Q_1$ is nonzero, hence isomorphic to the
standard form of degree 1, so Q is isomorphic to the standard form, which
proves 16.11.1.

It remains to prove the equivalence of (i) and (ii) in 16.9. We can
obviously suppose that S is the spectrum of an alg closed field, \var{\inserted[typed]{and we are}}{and we are}{insertion}{ega25_25_181-v0001} \var{\xout{on the other hand}}{}{typewriter strike-out}{ega25_25_181-v0002}
reduced to proving in this case \tikz[baseline=(lemma.base)]{\node[draw,ellipse,handwritten,text=black,inner sep=2pt] (lemma) {\shortstack[l]{\var{\inserted{le}}{le}{insertion}{ega25_25_181-v0003} \var{\inserted[typed]{\uline{Lemme} 16.11.2.}}{\uline{Lemme} 16.11.2.}{insertion}{ega25_25_181-v0004}\\\var{\deleted{\unsure{que}}}{}{deletion}{ega25_25_181-v0005} Q lisse \hand{$\Longleftrightarrow$} Q standart.}};}
% NOTE: The complete source diagram is retained verbatim, including its prose, under the exact-diagram translation protocol; its English wording is given in a translator note.
% NOTE: The pen oval encloses the lemma title and its formula; its upper edge crosses the preceding typed line without canceling that line.
By 16.2. we may suppose that k has char 2 and E has odd
rank. \var{\deleted{\var{\xout{We already know that (ii) implies that}}{}{typewriter strike-out}{ega25_25_181-v0007}}}{}{deletion}{ega25_25_181-v0006}
% NOTE: Tiny dark specks above the initial On and a small n-like remnant below it have no secure added-word or symbol reading.
 \var{\inserted[typed]{When}}{When}{insertion}{ega25_25_181-v0008} Q is isomorphic to the
standard form, \var{\xout{\unsure{one could}}}{}{typewriter strike-out}{ega25_25_181-v0009} one verifies at once, using 16.3. , that it
is smooth (the kernel of $\overset{\hand{\vee}}{E}$ \hand{$\longrightarrow$}Ev has dimension 1, and Q is not identically
zero on this kernel, so has only the trivial zero there); \var{\xout{therefore}}{}{typewriter strike-out}{ega25_25_181-v0010}
\var{\inserted[typed]{On the other hand}}{On the other hand}{insertion}{ega25_25_181-v0011} \var{\deleted{\var{\xout{(ii) \unsure{$\Rightarrow$} (i). Now}}{}{typewriter strike-out}{ega25_25_181-v0013}}}{}{deletion}{ega25_25_181-v0012} again by 16.3. , \var{\inserted[typed]{Q smooth}}{Q smooth}{insertion}{ega25_25_181-v0014} \var{\deleted{(ii)}}{}{deletion}{ega25_25_181-v0015} implies that \var{\xout{the kernel}}{}{typewriter strike-out}{ega25_25_181-v0016} the restriction
of Q to the kernel N of Q has only the trivial zero, which evidently
\var{\inserted[typed]{means}}{means}{insertion}{ega25_25_181-v0017} \var{\xout{implies}}{}{typewriter strike-out}{ega25_25_181-v0018} that N has dim 1 and that $Q|N \hand{\neq} 0$. Now taking a complement
of N, one sees at once that this implies that Q is isomorphic to the
standar\unsure{d|t} form\hand{/}. This \var{\inserted[typed]{proves 16.11.2. and}}{proves 16.11.2. and}{insertion}{ega25_25_181-v0019} completes the proof of 16.9.

NB We could have given 16.11.2. right at the beginning as a corollary to
16.2. or 16.3. \hand{(} which should have been interchanged); thus the part of the number
independent of the ``corrected discriminant'' would have been grouped at the beginning of the number.
Besides, you must know better than I do to what extent the notion
of corrected discriminant and 16.9. \unsure{were} known, so as to give credit where
it is due. Perhaps there is\unsure{-}already another accepted terminology? I note,
moreover, that it is the \uline{\unsure{cor}\unsure{rected}} discriminant that really deserves the name
of discriminant; it is this that generalizes to the discriminant of an
arbitrary form (cf the following number). The discriminant of a quadratic form
(in the terminology adopted here, and which, unless I am mistaken, is the

\mspage{\hand{181}}
\tnote{25-181-1}{The hand-encircled source diagram is retained literally. Its inserted title reads the Lemma 16.11.2; its line reads Q smooth if and only if Q standard, with a doubtful canceled that. The separate inserted article, typed lemma title, canceled word and handwritten ellipse remain distinct. Alg closed means algebraically closed, char means characteristic, and dim means dimension. The printed Ev following the arrow is retained.}
\tnote{25-181-2}{And retains the e/a doubt in et, were retains the e/t doubt in étaient, and the two parts of corrected retain the n/c and d/r doubts in corrigé. The doubtful hyphen in the French interrogative y a-t-il remains as an uncertain punctuation mark in the translated question. The author's request about priority and terminology is retained as a request.}
